\documentclass[11pt]{amsart}

% Encoding and fonts
\usepackage[T1]{fontenc}
\usepackage[utf8]{inputenc}
\usepackage{lmodern}

% Math packages
\usepackage{amsmath}
\usepackage{amssymb}
\usepackage{amsfonts}
\usepackage{amsthm}
\usepackage{mathtools}

% Layout and utilities
\usepackage{geometry}
\usepackage{enumitem}
\usepackage[hidelinks]{hyperref}

\geometry{margin=1in}

\numberwithin{equation}{section}
\setcounter{tocdepth}{3}
\setcounter{secnumdepth}{3}

% Theorem-like environments
\theoremstyle{plain}
\newtheorem{theorem}{Theorem}[section]
\newtheorem{proposition}[theorem]{Proposition}
\newtheorem{lemma}[theorem]{Lemma}
\newtheorem{corollary}[theorem]{Corollary}
\newtheorem{conjecture}[theorem]{Conjecture}

\theoremstyle{definition}
\newtheorem{definition}[theorem]{Definition}
\newtheorem{heuristic}[theorem]{Heuristic}
\newtheorem{experiment}[theorem]{Experiment}
\newtheorem{question}[theorem]{Question}

\theoremstyle{remark}
\newtheorem{remark}[theorem]{Remark}
\newtheorem{warning}[theorem]{Warning}

\title{Tubular Confinement and Clocked Chebyshev Cancellation}
\author{}
\date{}

\begin{document}

\maketitle

\tableofcontents

\section{Definitions}
\label{sec:definitions}

Let
\begin{equation}
\rho_k=\beta_k+i\gamma_k
\label{eq:zero_definition}
\end{equation}
denote a nontrivial zero of the Riemann zeta function, ordered by
increasing ordinate $\gamma_k$.

\begin{definition}[Logarithmic clock]
\label{def:logarithmic_clock}
For $k\ge 3$, define
\begin{equation}
\tau_k
=
\pi\frac{\log\log k}{\log k}.
\label{eq:tau_definition}
\end{equation}
\end{definition}

\begin{definition}[Free-amplitude phase-walk]
\label{def:free_amplitude_phase_walk}
For $\sigma\in\mathbb R$, define
\begin{equation}
z_k(\sigma)
=
\exp\big((\sigma+i\gamma_k)\tau_k\big),
\label{eq:zk_sigma_definition}
\end{equation}
and
\begin{equation}
S_N(\sigma)
=
\sum_{k=3}^{N}z_k(\sigma).
\label{eq:SN_sigma_definition}
\end{equation}
\end{definition}

\begin{definition}[Clocked Chebyshev scale]
\label{def:clocked_chebyshev_scale}
Define
\begin{equation}
x_k
=
\exp(\tau_k).
\label{eq:xk_definition}
\end{equation}
Equivalently,
\begin{equation}
\log x_k
=
\tau_k
=
\pi\frac{\log\log k}{\log k}.
\label{eq:log_xk_tau}
\end{equation}
\end{definition}

\begin{definition}[Tubular confinement]
\label{def:tubular_confinement}
The walk $S_N(\sigma)$ is said to be tubarly confined if there exist
$C\in\mathbb C$, $R>0$, and $N_0\ge 3$ such that
\begin{equation}
|S_N(\sigma)-C|\le R
\label{eq:tubular_confinement_definition}
\end{equation}
for every $N\ge N_0$.
\end{definition}

\begin{definition}[Clocked Chebyshev cancellation]
\label{def:clocked_chebyshev_cancellation}
The clocked Chebyshev terms are said to have controlled cancellation at
amplitude $\sigma$ if there exist $C\in\mathbb C$, $R>0$, and $N_0\ge 3$
such that
\begin{equation}
\left|
\sum_{k=3}^{N}
x_k^\sigma
\left(
\cos(\gamma_k\log x_k)
+
i\sin(\gamma_k\log x_k)
\right)
-C
\right|
\le R
\label{eq:clocked_chebyshev_cancellation_definition}
\end{equation}
for every $N\ge N_0$.
\end{definition}

\section{Equivalence theorem}
\label{sec:equivalence_theorem}

\begin{theorem}[Tubular confinement is clocked Chebyshev cancellation]
\label{thm:tube_chebyshev_equivalence}
For every $\sigma\in\mathbb R$, the following statements are equivalent.

\begin{enumerate}[label=\textnormal{(\roman*)}]
\item The walk $S_N(\sigma)$ is tubarly confined:
\begin{equation}
|S_N(\sigma)-C|\le R
\label{eq:theorem_tube_condition}
\end{equation}
for every $N\ge N_0$.

\item The clocked Chebyshev terms have controlled cancellation:
\begin{equation}
\left|
\sum_{k=3}^{N}
x_k^\sigma
\left(
\cos(\gamma_k\log x_k)
+
i\sin(\gamma_k\log x_k)
\right)
-C
\right|
\le R
\label{eq:theorem_chebyshev_condition}
\end{equation}
for every $N\ge N_0$.
\end{enumerate}
\end{theorem}

\begin{proof}
By \eqref{eq:xk_definition} and \eqref{eq:log_xk_tau}, one has
\begin{equation}
x_k^\sigma
=
\exp(\sigma\tau_k),
\label{eq:xk_sigma_expansion}
\end{equation}
and
\begin{equation}
\gamma_k\log x_k
=
\gamma_k\tau_k.
\label{eq:gamma_log_xk_identity}
\end{equation}
Therefore
\begin{align}
x_k^\sigma
\left(
\cos(\gamma_k\log x_k)
+
i\sin(\gamma_k\log x_k)
\right)
&=
\exp(\sigma\tau_k)
\left(
\cos(\gamma_k\tau_k)
+
i\sin(\gamma_k\tau_k)
\right)
\label{eq:clocked_term_expanded}
\\
&=
\exp(\sigma\tau_k)\exp(i\gamma_k\tau_k)
\label{eq:clocked_term_euler}
\\
&=
\exp\big((\sigma+i\gamma_k)\tau_k\big)
\label{eq:clocked_term_combined}
\\
&=
z_k(\sigma).
\label{eq:clocked_term_equals_zk}
\end{align}
Summing \eqref{eq:clocked_term_equals_zk} from $k=3$ to $N$ gives
\begin{equation}
\sum_{k=3}^{N}
x_k^\sigma
\left(
\cos(\gamma_k\log x_k)
+
i\sin(\gamma_k\log x_k)
\right)
=
S_N(\sigma).
\label{eq:partial_sum_identity}
\end{equation}

Assume first that $S_N(\sigma)$ is tubarly confined. Then
\eqref{eq:theorem_tube_condition} holds. By the identity
\eqref{eq:partial_sum_identity}, this is exactly
\eqref{eq:theorem_chebyshev_condition}. Thus tubular confinement implies
controlled clocked Chebyshev cancellation.

Conversely, assume that controlled clocked Chebyshev cancellation holds.
Then \eqref{eq:theorem_chebyshev_condition} holds. By
\eqref{eq:partial_sum_identity}, the partial sum appearing in
\eqref{eq:theorem_chebyshev_condition} is exactly $S_N(\sigma)$.
Therefore \eqref{eq:theorem_tube_condition} holds. Thus controlled
clocked Chebyshev cancellation implies tubular confinement.

The two statements are equivalent.
\end{proof}

\section{Critical amplitude}
\label{sec:critical_amplitude}

\begin{proposition}[Critical amplitude]
\label{prop:critical_amplitude}
At
\begin{equation}
\sigma=\frac12,
\label{eq:sigma_half}
\end{equation}
the clocked Chebyshev amplitude is
\begin{equation}
x_k^\sigma=x_k^{1/2}.
\label{eq:critical_amplitude_identity}
\end{equation}
Thus $\sigma=1/2$ is the critical amplitude level of the clocked
Chebyshev model.
\end{proposition}

\begin{proof}
Substitution of \eqref{eq:sigma_half} into \eqref{eq:xk_sigma_expansion}
gives
\begin{equation}
x_k^\sigma=x_k^{1/2}.
\label{eq:critical_amplitude_substitution}
\end{equation}
For a zero $\rho=\beta+i\gamma$, the standard zero term has the
algebraic form
\begin{equation}
x^\rho
=
x^\beta
\left(
\cos(\gamma\log x)
+
i\sin(\gamma\log x)
\right).
\label{eq:xrho_decomposition}
\end{equation}
Hence the exponent $\beta=1/2$ corresponds to the amplitude $x^{1/2}$.
In the clocked model, the free parameter $\sigma$ plays the same formal
amplitude role, and $\sigma=1/2$ is the critical value.
\end{proof}

\section{Conclusion}
\label{sec:conclusion}

The exact identity
\begin{equation}
S_N(\sigma)
=
\sum_{k=3}^{N}
x_k^\sigma
\left(
\cos(\gamma_k\log x_k)
+
i\sin(\gamma_k\log x_k)
\right)
\label{eq:conclusion_identity}
\end{equation}
with
\begin{equation}
\log x_k
=
\pi\frac{\log\log k}{\log k}
\label{eq:conclusion_clock}
\end{equation}
shows that the logarithmic phase-walk is not merely analogous to a
Chebyshev-type oscillatory sum. It is exactly a Chebyshev-shaped
oscillatory sum evaluated on the clocked scale \eqref{eq:conclusion_clock}.

Therefore,
\begin{equation}
\boxed{
\text{tubular confinement of }S_N(\sigma)
\Longleftrightarrow
\text{controlled cancellation of clocked Chebyshev-type terms}
}
\label{eq:final_equivalence}
\end{equation}
for every real $\sigma$.

The critical case is
\begin{equation}
\boxed{
\sigma=\frac12,
}
\label{eq:final_critical_sigma}
\end{equation}
because this is exactly the value for which the clocked amplitude is
$x_k^{1/2}$.

\end{document}